\subsection*{一种对相同比例变化给相同读数的刻度}
从一百增加到一百零一，与从一千增加到一千零一，绝对变化相同，但相对变化不同。自然对数在局部测量的是 $dx/x$，因此更接近按比例的变化尺度。相同百分比增幅产生相同对数增量。

因为 $\ln$ 与指数互逆，对 $\mathrm e^{\ln x}=x$ 求导，得到 $x(\ln x)'=1$，所以导数为 $1/x$。定义域要求 $x>0$，分母也说明不能在零点套用这个推导。

\subsection*{用面积把乘法法则重新看一遍}
曲线 $1/t$ 从一到 $x$ 的有向面积等于 $\ln x$。区间同比例缩放时，高度相应反比例变化，面积保持；这为乘法转换成加法提供了另一种几何解释。

当 $x$ 从一百变为一百零一，对数增量为 $\ln(1.01)$，一阶近似是 $0.01$。但若从一变为二，直接把对数增量当一就不准确，因为变化已不够小。局部相对变化近似不能当成全局等式。

读对数图时也要区分“增长量下降”和“相对增长率下降”。指数曲线在对数坐标中成为直线，因为它每单位时间增加相同的对数值。这个表示很适合检查比例增长假设，但图形近似直线仍不足以证明具体机制。
